% F5c -- The other branch: one prompt sent to a frozen model.
\begin{figure}[t]
\centering

{\footnotesize\textbf{One of the prompts of the prompting branch,} as it is sent
to a frozen model. Quoted from Tiphain and Benjamin's work and shortened at the
two ellipses; nothing else is altered.\par}
\vspace{3pt}

\begin{lstlisting}
User's day started at 00:00 at station VSOYRO. [...] They stayed at BSOROT
until 05:34, then left for station VOYQZ, arriving at 06:15. [...] It is now
07:45 and the user is at VSOYRO.
Predict the station the user will visit next. Rely only on their movement
pattern and the stations around them. Possible stations include: BKVOR4,
BKVOF2, BKVS2D, BSOROT, BSVOLOK, DRVOAN, DSXVOP, LKVTSE. Which station are
they most likely to go to next?
Look at the user's daily flow (loops, pauses, commuting, travel habits).
Explain your choice in a few natural sentences, then end with the station in
this format:
FINAL ANSWER: [STATION]
\end{lstlisting}

\caption{What one prompt looks like}
\label{fig:prompting}
\figtext{%
The same problem, asked the other way. Three differences with the training
record of Figure~\ref{fig:training-example} are worth seeing, because they are
what makes the two branches answer different questions. The day is narrated in
English, stop by stop and with clock times, so one event costs a whole clause
where the fine-tuned model reads two symbols, and a long user stops fitting in
the model's window long before a short one does. The candidate cells are handed
over inside the prompt, eight of them here, and the answer comes back as free
text that has to be parsed, which is what the \texttt{FINAL ANSWER} convention
is for; the fine-tuned model instead chooses among 369 symbols it already
carries in its vocabulary. And since no weight ever moves, nothing learned
about one user is available for the next: everything the model is to know about
this population has to be restated in every prompt, whereas the whole point of
the branch this report takes is that the population is what the weights end up
holding. Neither is a better instrument. Asking a general model what it already
knows about a day described in words, and asking what can be learned from the
population itself, are two questions, and the second is this report's.}
\end{figure}
